% !PRINTED 134
\weberpage{134}
%<<CONT
\weq{5}{
f_2\!\left(\frac{\gamma+\delta\omega}{\alpha+\beta\omega}\right) = \left(\frac{2}{\beta}\right)\varrho\, e^{\frac{3\pi i}{8}\beta(2\delta-\alpha)}\, f_1(\omega), \qquad \alpha \equiv 0 \pmod{2},
}
und ersetzt man hierin $\omega$ durch $\omega+3$ und $\gamma, \alpha$ durch $\gamma-3\delta, \alpha-3\beta$, so folgt
\weq{6}{
f_2\!\left(\frac{\gamma+\delta\omega}{\alpha+\beta\omega}\right) = -\varrho\, e^{\frac{3\pi i}{8}\beta(2\delta-\alpha)}\, f(\omega), \qquad \alpha-\beta \equiv 0 \pmod{2}.
}

Setzt man in (4), (5), (6)
\[
f_2\!\left(\frac{\gamma+\delta\omega}{\alpha+\beta\omega}\right) = f_1\!\left(-\frac{\alpha+\beta\omega}{\gamma+\delta\omega}\right),
\]
und vertauscht dann $\alpha, \beta, \gamma, \delta$ mit $-\gamma, -\delta, \alpha, \beta$, so ergibt sich
\weq{7}{
f_1\!\left(\frac{\gamma+\delta\omega}{\alpha+\beta\omega}\right) = \left(\frac{2}{\gamma}\right)\varrho\, e^{-\frac{3\pi i}{8}\gamma(2\alpha-\delta)}\, f_2(\omega), \qquad \delta \equiv 0 \pmod{2},
}
\weq{8}{
{}= \left(\frac{2}{\delta}\right)\varrho\, e^{-\frac{3\pi i}{8}\delta(2\beta+\gamma)}\, f_1(\omega), \qquad \gamma \equiv 0 \pmod{2},
}
\weq{9}{
{}= -\varrho\, e^{-\frac{3\pi i}{8}\delta(2\beta+\gamma)}\, f(\omega), \qquad \gamma-\delta \equiv 0 \pmod{2}.
}

Aus (9) und (6) erhält man, indem man $\omega$ durch $\dfrac{-\gamma+\alpha\omega}{\delta-\beta\omega}$ und dann $\alpha, \beta, \gamma, \delta$ durch $\delta, -\beta, -\gamma, \alpha$ ersetzt:
\weq{10}{
f\!\left(\frac{\gamma+\delta\omega}{\alpha+\beta\omega}\right) = -\varrho\, e^{-\frac{3\pi i}{8}\alpha(2\beta+\gamma)}\, f_1(\omega), \qquad \alpha-\gamma \equiv 0 \pmod{2},
}
\weq{11}{
f\!\left(\frac{\gamma+\delta\omega}{\alpha+\beta\omega}\right) = -\varrho\, e^{\frac{3\pi i}{8}\beta(2\alpha-\delta)}\, f_2(\omega), \qquad \beta-\delta \equiv 0 \pmod{2},
}
und wenn man endlich in (10) $\omega$ durch $\omega+9$ und $\gamma, \alpha$ durch $\gamma-9\delta, \alpha-9\beta$ ersetzt:
\weq{12}{
f\!\left(\frac{\gamma+\delta\omega}{\alpha+\beta\omega}\right) = \varrho\left(\frac{2}{\alpha-\beta}\right) e^{-\frac{3\pi i}{8}(\alpha-\beta)(\alpha+\beta+\gamma-\delta)}\, f(\omega),
}
\[
\alpha+\beta+\gamma-\delta \equiv 0 \pmod{2}.
\]
\footnote{Die von \gesperrt{Hermite} (Sur la théorie des équations modulaires, Paris 1859) eingeführten Funktionen $\varphi(\omega), \psi(\omega), \chi(\omega)$ hängen mit den Funktionen $f(\omega), f_1(\omega), f_2(\omega)$ durch die Gleichungen
\[
f(\omega) = \frac{\sqrt[6]{2}}{\chi(\omega)}, \quad f_1(\omega) = \sqrt[6]{2}\,\frac{\psi(\omega)}{\chi(\omega)}, \quad f_2(\omega) = \sqrt[6]{2}\,\frac{\varphi(\omega)}{\chi(\omega)}
\]
zusammen. Die Transformationsformeln der $f$-Funktionen lassen sich mit Benutzung der Relation $\alpha\delta-\beta\gamma=1$ auf mannigfaltige Weise umgestalten. So sind die von \gesperrt{Hermite} a.\,a.\,O.\ angegebenen Formeln nicht ohne weiteres mit den unserigen als identisch zu erkennen. Eine einfache Rechnung zeigt aber ihre Übereinstimmung. Die oben gegebenen Formeln haben den Vorzug, daß sie, ohne an Einfachheit zu verlieren, je zwei der sechs Transformationsklassen in einen Ausdruck zusammenfassen.}

\begin{center}\rule{2cm}{0.4pt}\end{center}
